\chapter{Fazit}
\label{sec:fazit}

\section{Auswertung und Zusammenfassung}

Die im Rahmen der Projektarbeit definierten Zielsetzungen wurden erfüllt. Die Ergebnisse zeigen, dass es mit handelsüblicher Hardware (USB-Webcam) und sehr kostengünstiger Software möglich ist, bereits verwertbare Ergebnisse zu generieren. Ursprünglich war angedacht, eine Open-Source-Lösung als Fundament für den Eye Tracker zu verwenden. Nach ausgiebiger Marktrecherche konnte ein passendes Softwarefundament ausfindig gemacht werden. Der Beam Eye Tracker ist zwar nicht Open Source, aber die Daten, die er durch seine API bereitstellen kann, bieten genug Spielraum, um individuelle Lösungen damit bauen zu können. Eine Verschmelzung zwischen CARLA und dem Eye Tracker ließ sich ohne signifikante technische Herausforderungen realisieren. Die Daten aus dem Instance Segmentation Sensor können zusammen mit den Werten des Eye Trackers ausgewertet werden. Somit ist es möglich, nicht nur die geschätzten Koordinaten des Eye Trackers auf der Leinwand auszugeben, sondern auch das betrachtete Objekt innerhalb der CARLA Simulationsumgebung zu bestimmen.
Während der Implementierung des Eye Tracking Systems sind an manchen Stellen auch Probleme aufgetreten. Die Belichtung war vor allem während der Kalibrierung ein Problem. Durch verschiedene Ansätze und Experimente im Fahrsimulator konnten hier mit externer Belichtung gute Ergebnisse erzielt werden.
Eine weitere Herausforderung stellten die durch die gekrümmte Leinwand verfälschten Tracking Daten dar. Durch umfassende Analyse des Problems, einige mathematische Ansätze und Experimente im Fahrsimulator, konnte die daraus resultierende Abweichung signifikant reduziert werden. Die in Abbildung \ref{fig:imag12} dargestellten Ergebnisse zeigen, dass die korrigierten Daten eine deutlich verbesserte Genauigkeit aufweisen.
Die Tracking Daten sind allerdings nicht perfekt und lassen sich durch äußere Einflüsse wie schlechte Belichtung, Bedienungsfehler bei der Kalibrierung, Veränderung der Sitzposition nach der Kalibrierung oder zwei Personen gleichzeitig im Fahrsimulator negativ beeinflussen. Der Eye Tracker könnte allerdings zum aktuellen Entwicklungsstand bereits funktional im Fahrsimulator eingesetzt werden. Anhand der gemessenen Performance, deuten die Ergebnisse darauf hin, dass bereits mit dem aktuellen Entwicklungsstand verwertbare Blickdaten generiert werden können. Im Ausblick sollen weitere Schritte und mögliche Ideen präsentiert werden, mit welchen über diesen Low Budget Ansatz hinaus noch bessere Ergebnisse erzeugt werden könnten.


\section{Ausblick}

Aufbauend auf den Ergebnissen dieser Arbeit kann eine weitere Verbesserung des Testaufbaus erfolgen. Dazu zählen eine Ausbesserung der Lichtverhältnisse im Fahrsimulator, eine fest verbaute Kamera und gegebenenfalls eine Blende, sodass nur der Fahrer und nicht zusätzlich der Beifahrer von der Kamera erkannt wird. Ebenso kann die Eye Tracker Genauigkeit verbessert werden, indem die Korrektur für den gekrümmten Bildschirm durch einen potentiell noch besseren Optimierungsalgorithmus durchgeführt wird. Eine weitere Optimierungsmöglichkeit stellt die Anpassung der vertikalen Kalibrierung dar. Bspw. ließe sich dadurch ein statisches Offset korrigieren, wenn die Person während einer Fahrt in vertikale Richtung die Sitzposition verändert. Eine Integration von künstlicher Intelligenz zur Bestimmung der genauen Blickposition könnte ebenfalls die Genauigkeit verbessern. 
Abschließend bieten sich weiterführende Analysen an, wie eine zeitabhängige Auswertung von Blickpositionen in einem Zeitstrahl. 
Der Code des Eye Trackers ist dem Anhang beigefügt und kann somit weiterentwickelt werden.


\clearpage
